% Adapted from Jake's Resume (MIT License)
\documentclass[letterpaper,10pt]{article}

\usepackage{cmap}
\usepackage[T2A,T1]{fontenc}
\usepackage[utf8]{inputenc}
\usepackage[english,russian]{babel}
\usepackage[empty]{fullpage}
\usepackage{titlesec}
\usepackage{enumitem}
\usepackage[hidelinks]{hyperref}
\usepackage{fancyhdr}
\usepackage{tabularx}
\usepackage{xcolor}
\usepackage{microtype}
\input{glyphtounicode}
\pdfgentounicode=1

\pagestyle{fancy}
\fancyhf{}
\renewcommand{\headrulewidth}{0pt}
\renewcommand{\footrulewidth}{0pt}

\addtolength{\oddsidemargin}{-0.62in}
\addtolength{\evensidemargin}{-0.62in}
\addtolength{\textwidth}{1.24in}
\addtolength{\topmargin}{-0.66in}
\addtolength{\textheight}{1.30in}

\urlstyle{same}
\raggedbottom
\raggedright
\setlength{\tabcolsep}{0in}
\setlength{\parindent}{0pt}

\titleformat{\section}{\vspace{-3pt}\raggedright\large}{}{0em}{}[\color{black}\titlerule \vspace{-5pt}]

\newcommand{\resumeItem}[1]{\item\small{#1\vspace{-1pt}}}
\newcommand{\resumeSubheading}[4]{
  \vspace{0pt}\item
  \begin{tabular*}{0.98\textwidth}[t]{l@{\extracolsep{\fill}}r}
    \textbf{#1} & #2 \\
    \textit{\small #3} & \textit{\small #4} \\
  \end{tabular*}\vspace{-6pt}
}
\newcommand{\resumeSubSubheading}[2]{
  \item
  \begin{tabular*}{0.98\textwidth}{l@{\extracolsep{\fill}}r}
    \textit{\small #1} & \textit{\small #2} \\
  \end{tabular*}\vspace{-6pt}
}
\newcommand{\resumeProjectHeading}[2]{
  \item
  \begin{tabular*}{0.98\textwidth}{l@{\extracolsep{\fill}}r}
    \small #1 & \small #2 \\
  \end{tabular*}\vspace{-6pt}
}
\renewcommand\labelitemii{$\vcenter{\hbox{\tiny$\bullet$}}$}
\newcommand{\resumeSubHeadingListStart}{\begin{itemize}[leftmargin=0.14in,label={}]}
\newcommand{\resumeSubHeadingListEnd}{\end{itemize}}
\newcommand{\resumeItemListStart}{\begin{itemize}[leftmargin=0.22in,itemsep=1pt,topsep=2pt]}
\newcommand{\resumeItemListEnd}{\end{itemize}\vspace{-5pt}}

\hypersetup{
  pdftitle={Арсений Миронов - Резюме},
  pdfauthor={Арсений Миронов},
  pdfsubject={Резюме Machine Learning и LLM Inference Engineer}
}

\begin{document}

\begin{center}
  \textbf{\Huge Арсений Миронов} \\ \vspace{2pt}
  \small Нижний Новгород, Россия $|$ +7 (917) 618-94-15 $|$
  \href{mailto:arseniy.mironov.dev@gmail.com}{\underline{arseniy.mironov.dev@gmail.com}} \\
  \href{https://github.com/Napkin-AI}{\underline{github.com/Napkin-AI}} $|$
  \href{https://t.me/xNapkin}{\underline{t.me/xNapkin}}
\end{center}

\section{Образование}
\resumeSubHeadingListStart
  \resumeSubheading
    {ННГУ им. Н. И. Лобачевского}{Нижний Новгород, Россия}
    {Магистратура. Прикладная математика и информатика}{2026 -- н. в.}
  \resumeSubSubheading
    {Бакалавриат. Фундаментальная информатика и информационные технологии}{2022 --- 2026}
  \resumeSubheading
    {Школа анализа данных Яндекса}{ }
    {Направление Data Science}{2 курс}
\resumeSubHeadingListEnd

\section{Опыт работы}
\resumeSubHeadingListStart
  \resumeSubheading
    {Инженер}{Янв. 2026 -- н. в.}
    {Huawei}{Нижний Новгород, Россия}
  \resumeItemListStart
    \resumeItem{Развиваю стек инференса SGLang для Ascend NPU с фокусом на инференс LLM и diffusion-моделей, attention backends и оптимизацию производительности.}
    \resumeItem{Интегрировал sparse attention backends в SGLang Diffusion для Ascend NPU, включая конфигурацию backend'ов, проверку точности, документацию и профилирование производительности.}
    \resumeItem{Реализовал поддержку mxfp8 flash attention backend, квантованных моделей Wan и Flux (w8a8f8/w4a4f8).}
    \resumeItem{Исследовал ускорение prefill для Qwen3 и diffusion-LLM, а также узкие места инференса на Ascend NPU.}
    \resumeItem{Список всех PRs: \href{https://github.com/sgl-project/sglang/pulls?q=is\%3Apr+state\%3Aclosed+author\%3ANapkin-AI}{\underline{sgl-project/sglang/pulls}}.}
  \resumeItemListEnd

  \resumeSubheading
    {Инженер-стажёр}{Июль 2025 -- Авг. 2025}
    {Huawei}{Нижний Новгород, Россия}
  \resumeItemListStart
    \resumeItem{Исследовал раздельный запуск этапов prefill и decode в vLLM на аппаратных ускорителях.}
    \resumeItem{Адаптировал KV-cache transfer engine под архитектуру целевого ускорителя и провёл бенчмарки производительности LLM при disaggregated serving.}
  \resumeItemListEnd
\resumeSubHeadingListEnd

\section{Open Source и проекты}
\resumeSubHeadingListStart
  \resumeProjectHeading
    {\textbf{OpenCV --- оптимизация детектирования QR и ArUco} $|$ \emph{C++}}{}
  \resumeItemListStart
    \resumeItem{Реализовал \texttt{cv::approxPolyN} --- алгоритм жадной аппроксимации контура выпуклым многоугольником с настраиваемым числом сторон и порогом ошибки по площади.}
    \resumeItem{Реализация \href{https://github.com/opencv/opencv/pull/25607}{\underline{Pull Request \#25607}} вошла в OpenCV 4.11.0.}
  \resumeItemListEnd
\resumeSubHeadingListEnd

\section{Публикации и участие в конференциях}
\resumeSubHeadingListStart
  \resumeProjectHeading
    {\textbf{А. И. Миронов}, Н. А. Соколов, А. А. Серебрякова}{2025}
  \resumeItemListStart
    \resumeItem{«Влияние объема синтетически сгенерированных изображений в наборе данных для задачи сегментации аксонов в электронной микроскопии мозга». \textit{GraphiCon 2025}. \href{https://doi.org/10.25686/978-5-8158-2474-4-2025-828-834}{\underline{DOI}}}
  \resumeItemListEnd
\resumeSubHeadingListEnd

\section{Спортивное программирование}
\resumeSubHeadingListStart
  \resumeProjectHeading{\textbf{ICPC 2025--2026}}{}
  \resumeItemListStart
    \resumeItem{Квалификация в Нижнем Новгороде: диплом 2 степени; четвертьфинал в Саратове: диплом 3 степени; полуфинал ICPC, финал NERC в Санкт-Петербурге: диплом 3 степени.}
  \resumeItemListEnd
  \resumeProjectHeading{\textbf{ICPC 2024--2025}}{}
  \resumeItemListStart
    \resumeItem{Квалификация в Нижнем Новгороде: диплом 2 степени; четвертьфинал в Саратове: 59 место.}
  \resumeItemListEnd
\resumeSubHeadingListEnd

\section{Технические навыки}
\begin{itemize}[leftmargin=0.14in,label={}]
  \small{\item{
    \textbf{Языки:} Python, C++, SQL, Triton \\
    \textbf{ML / Inference:} PyTorch, SGLang, vLLM, OpenCV, NumPy, pandas, scikit-learn, Matplotlib \\
    \textbf{Системы / инструменты:} Linux, Git, Jupyter, LaTeX \\
    \textbf{Английский:} B1
  }}
\end{itemize}

\end{document}
